\newpage
\section{ANNEX 4 - The EURO-LABS Legacy: an ambitious successful experiment in scientific collaboration} 
\subsection{Overview}
Building on more than two decades of European experience with Transnational Access, EURO-LABS brought together nuclear physics, and accelerator science and detector technology for particle physics within a single coordinated European programme. EURO-LABS built on this successful experience while testing a more ambitious model. EURO-LABS, was Europe’s first experiment of amalgamating different communities.  European Nuclear and Particle Physics continues to be in an evolutionary stage and thus requires a continuous scientific and technical growth and thus requiring the best  the use of Europe's distributed accelerator infrastructure. By coordinating access to complementary facilities, rather than treating them as isolated resources, and leveraging expertise gained in the field curiosity-driven research and thus address fundamental and technical questions were able to be performed  .  By  optimizing the use of Europe's distributed accelerator infrastructure by coordinating access to complementary facilities rather than treating them as isolated resources. The most appropriate facility for a scientific question is not necessarily the largest facility, or the facility closest to the user, but the facility with the capabilities best suited to the research. Large, medium and smaller facilities each contributed distinctive capabilities and expertise to the European ecosystem. This approach allowed the project to exploit the diversity of European infrastructures, including facilities of different sizes and with different areas of expertise. Their efficient use was the key element in this endeavour, that maximizes the investment put in Europe to maintain its leadership in this field  The project also demonstrated that complementary research infrastructures, technologies and expertise of these communities can generate value well beyond the capabilities of individual facilities or communities. The important result is therefore not simply the quantity of access, but the ability to match the right scientific question with the most appropriate European capability.

By coupling access to continuous improvement of facilities and services, EURO-LABS created a feedback loop in which today's users contribute to the capabilities available to tomorrow's users. The project recognised that the value of access does not end when an experiment is completed. Scientific needs identified by users can lead to improvements in instrumentation, detectors, accelerator operation, data acquisition, software and other services. Service Improvements therefore became an integral component of the access model rather than a separate engineering activity. In addition to the heart of this project namely Transnational Access for science and technology, it was intelligently combined among others related activities. These included critical hands-on training of the next  generation. This approach is particularly valuable for early-career researchers because it develops not only knowledge but practical experimental capability and professional networks. Complementary theoretical support for better ideas, interpretation of data and allied issues related to open data practices and the use of AI for more efficient operations and applications to e.g., medicine were some of the related issues that we treated.

In a nutshell its distinctive approach combined:
\begin{itemize}
\item coordinated Transnational Access to Europe's diverse accelerator facilities; 
\item Service Improvements driven by scientific and user needs; 
\item Joint Research Activities and technology development; 
\item extensive hands-on training for the next generation; 
\item theoretical support for experiment design and interpretation; 
\item Open Science and FAIR data practices
\item machine learning and AI for more efficient operation and optimisation; 
\item and pathways towards applications in areas including medicine, energy, materials and industry. 
\end{itemize}

EURO-LABS has explored the limits of how diverse communities could be having synergies and run a successful Transnational Access programme. The follow-up HORIZON EUROPE flagship Transnational Access programme within the INFRA-SERV call is set to include 6 fields within a pilot project, a clear recognition of the success of EURO-LABS. The EURO-LABS team, through coordinated, coherent and synergetic approach to science and technologies and by optimizing finances, showed the project of combining three communities can lead to a grand success. The project has answered the goals of EU in the work program at nominal costs by taking the strategy of minimizing the payment for the beam hours delivered by various Research Infrastructures (RI). Thus the project was also able to fund a larger number of important related activities as mentioned earlier. All the activities of the Transnational Access programme were reviewed by external experts. The Transnational Access provided was way beyond (almost a factor of $\sim$ 1.5) to what has been pledged to exploit all the relevant facilities over Europe. EURO-LABS gave the opportunity to start small, create new ideas and become big and innovative (including theory support). The optimum and coordinated use the eco system of variety of RI and their various facilities enhanced each other and created synergies. The project very strongly supported a diverse user community with an emphasis on new users and early career scientists. The scientific and technical excellence can be seen through high impact publication and presentations at conferences increasing the impact of basic science but also allied fields and technical achievements. In selected cases, technologies developed through the research ecosystem progressed towards transfer to hospitals and industrial partners.

EURO-LABS did not attempt to merge distinct scientific programmes in nuclear and particle physics. Rather, it demonstrated that substantial synergies exist in the capabilities that underpin those programmes. These include:
\begin{itemize}
\item accelerators and beam instrumentation; 
\item detectors and sensors; 
\item electronics and data acquisition; 
\item targets and instrumentation; 
\item simulation and software; 
\item machine learning and AI; 
\item data management and Open Science; 
\item theory–experiment interaction. 
\end{itemize}
This provided an important model for future integration: different scientific communities can retain their distinct missions while sharing the technologies, expertise and infrastructures that enable them.

Periodic Report 3 covers the final reporting period of the project. This appendix has a different purpose  which is provide one place very briefly the scientific, technical and other highlights of   EURO-LABS over the full four years

\subsection{RIs Landscape}
EURO-LABS started in September 2022 and will run until August 2026. It brings together 33 partner institutions from 18 European and non-European countries, involving 45 leading RI facilities (Fig.~\ref{fig:facilities}).
\begin{figure}[!htb]
   \centering
   \includegraphics[width=0.75\columnwidth]{MOP7145_f1.pdf}
    \caption{Map of EURO-LABS partner institutions and participating RIs across Europe. Associated-partner facilities in the US and Japan, are not shown here.}
   \label{fig:facilities}
\end{figure}
%\vspace{-0.3mm}
\subsubsection{Nuclear Physics RIs}

The Nuclear Physics RIs within EURO-LABS provide European and international users with TNA access to world-leading facilities delivering high-intensity stable ion beams, in some cases reaching the p$\mu$A range, covering isotopes from hydrogen to uranium. The broad EURO-LABS portofolio extends to facilities producing high-energy radioactive ion beams, nuclei not existing on Earth, using in-flight fragmentation techniques (of nano-sec lifetime). Post-accelerated radioactive ion beams (of msec half-time) using ISOL (Isotope Separation On-Line) methods offered for TNA access, as well as to neutron-beam facilities. 

For the first time, EURO-LABS also introduced a \textbf{Virtual Access (VA)} scheme offering access to the European Centre for Theoretical Studies in Nuclear Physics and Related Areas, and to theoretical support for experiments. Table~\ref{tabl:wp2-ris} summarizes the participating TNA and VA facilities.
% \vspace{-0.7mm}
\begin{table}[htb]
\centering
\caption{Participating Facilities in Nuclear Physics. ESFRI roadmap RIs are shown in bold}
\small
\begin{tabular}{
p{0.14\columnwidth}
p{0.16\columnwidth}
p{0.24\columnwidth}
p{0.24\columnwidth}
}
\toprule
\textbf{Type} & \textbf{Country} & \textbf{Host Lab} & \textbf{Infrastructure} \\
\midrule
S,R,N & Italy & INFN & LNL-NSDBF; LNS-AIPF \\
S,R,N & France & GANIL & \textbf{SPIRAL2} \\
S,R,N & France & CNRS-IJCLab & ALTO \\
S,R   & Germany & GSI & \textbf{FAIR} \\
S,R   & Finland & Univ.\ Jyvaskula & JYFL \\
R     & Switzerland & CERN & ISOLDE \\
N     & Switzerland & CERN & nTOF \\
S     & Poland & Univ.\ Warsaw & NLC-SLCJ \\
S     & Poland & IFJPAN & NLC-CCB \\
S     & Romania & IFIN-HH & Tandem \\
S,N   & Spain & USE & \rdelim\}{3}{1mm}[CLEAR] \\
S,N   & Hungary & ATOMKI &  \\
S,N   & Portugal & IST &  \\
Theory & Italy & ECT & ECT \\ \hline
Theory/VA & Poland & IFJPAN & Meanfield4Exp \\
Theory/VA & Spain & USE & Reaction4Exp \\
Theory/VA & Italy & U.Milano & Structure4Expt \\
\bottomrule
\end{tabular}

\vspace{1mm}
\raggedright
\footnotesize
S = Stable ion beams, R = Radioactive ion beams, N = Neutron beams.
\label{tabl:wp2-ris}
\end{table}
The scientific goal of the NP Transnational Access is to provide enhanced opportunities for exploring nuclei under extreme conditions: at high temperature ($T$), high angular momentum ($L$), large isospin ($N/Z$), and extreme masses ($A$), including nuclei close to the proton and neutron driplines. The supported TNA projects focus on addressing key open questions in the field, including among others the emergence of nuclear structure from nucleon interactions and the limits of nuclear stability and production of new chemical elements. The programme also supports innovative applications, including developments related to ion-beam cancer therapy.

\subsubsection{Accelerator R\&D RIs}

This activity provides TNA  to a broad range of unique and specialized accelerator R\&D facilities supporting the development and validation of concepts for future accelerator infrastructures, including studies relevant for next-generation colliders such as the CERN Future Circular Collider (FCC) and the Muon Collider. The participating facilities enable experimental investigations of advanced accelerator technologies, including superconducting and normal-conducting radio-frequency cavities, magnets, novel accelerator materials, and advanced acceleration schemes. Table~\ref{tabl:wp3-ris} summarizes the portfolio of participating facilities.
\begin{table}[htb]
\centering
\caption{Participating facilities in Accelerator R\&D. ESFRI roadmap RIs are shown in bold}
\small
\begin{tabular}{
p{0.12\columnwidth}
p{0.16\columnwidth}
p{0.24\columnwidth}
p{0.28\columnwidth}
}
\toprule
\textbf{Type} & \textbf{Country} & \textbf{Host Lab} & \textbf{Infrastructure} \\
\midrule
MT & Switzerland & CERN & HiRadMat \\
TI & Sweden & Univ.\ Uppsala & FREIA \\
TI & Italy & INFN-MI & LASA \\
TI & Italy & INFN-Univ.\ Sal & THOR \\
TI & France & CNRS-JCLab & SUPRATECH \\
TI & France & CEA Irfu-Syn\-ergium & MACHAFILM; CRYOMECH \\
TI & Switzerland & CERN & XBOX \\
EP & Germany & KIT-ALFA & KARA; FLUTE \\
EP & UK & STFC & CLARA, VELA \\
EP & Italy & INFN-Frascati & BTF; \textbf{SPARClab} \\
EP & France & CEA LiDyL & LPA-UHI100 \\
APP & Poland & INTC & RAPID \\
APP & Switzerland & CERN & CLEAR \\
\bottomrule
\end{tabular}

\vspace{1mm}
\raggedright
\footnotesize
MT = Material Testing,
TI = Technology Infrastructure,
EP = Electron \& Plasma Beam,
APP = Applications.
\label{tabl:wp3-ris}
%\vspace{-3mm}
\end{table}

\subsubsection{HEP Detector R\&D RIs}

The RIs included in this area provide TNA access to a broad range of facilities delivering energetic particle beams (protons, mesons, muons, and electrons), as well as irradiation facilities operating with mixed hadron and $\gamma$ fields. These facilities, summarized in Table~\ref{tabl:wp4-ris}, support detector prototype testing and characterisation studies. Such measurements are essential for understanding the behaviour and radiation tolerance of detectors and associated instrumentation under realistic beam conditions, supporting detector upgrades for HL-LHC operation and the development of future detector technologies.

\begin{table}[htb]
\centering
\caption{Participating Facilities in HEP Detector R\&D.ESFRI roadmap RIs are shown in bold}
\small
\begin{tabular}{
p{0.12\columnwidth}
p{0.16\columnwidth}
p{0.24\columnwidth}
p{0.26\columnwidth}
}
\toprule
\textbf{Type} & \textbf{Country} & \textbf{Host Lab} & \textbf{Infrastructure} \\
\midrule
TB & Switzerland & CERN & \textbf{PS} \& \textbf{SPS} \\
TB & Germany & DESY & DESY-II \\
DC & Croatia & RBI & RBI-AF \\
DC & Spain & ITAINNOVA & EMCLab \\
DC & Switzerland & CERN & IRRAD; GIF++ \\
DC & Slovenia & JSI & TRIGA Reactor \\
DC & Poland & IFJPAN & AIC-144 \\
DC & Belgium & UCLouvain & CRC \\
DC & UK & UoB & MC40 Cyclotron \\
\bottomrule
\end{tabular}

\vspace{1mm}
\raggedright
\footnotesize
TB = Test Beams,
DC = Detector Characterization.
\label{tabl:wp4-ris}
\end{table}

%\vspace{-1mm}
EURO-LABS has been running successfully since its start on September 1st, 2022. Information on the facilities and application templates were provided at the project website~\cite{euro-labs}. Applications included a scientific proposal describing the scope of the project, with particular emphasis on excellence and novelty, together with the list of team members requesting support. The TNA programme provided funding covering experimental costs as well as travel support for participating users. 
Applications were collected at the facility level, evaluated by a \textbf{User Selection Panel (USP)} composed of experts in the relevant scientific domain and facility area, with the final approval endorsed by the corresponding facility scientific and planning committee. 

EURO-LABS has been widely promoted through presentations at conferences, scientific meetings, and outreach activities within the related research communities. By \textbf{December 2025, a total of 490 TNA projects had been approved, supporting 3557 users across all participating facilities.}

%Beyond the Transnational Access programme, EURO-LABS also supported targeted service improvements at several participating %facilities, aimed at enhancing the available performance, operational capabilities, and parameter range of the RIs for the %benefit of the user community. These activities included the development of novel beam-control approaches based on 
%\textbf{machine-learning (ML)} techniques to improve beam quality, transport efficiency, and operational %reproducibility~\cite{geoff}.

%In Nuclear Physics, significant effort was additionally devoted to integrating community data and services within the %European Open Science Cloud (EOSC) framework, with the objective of strengthening the implementation of FAIR (Findable, %Accessible, Interoperable, and Reusable) data principles across the participating infrastructures.

\subsection{EURO-LABS: Science and Technology Advancing Together}
Over four years, EURO-LABS has demonstrated the impact that can be achieved by bringing together Europe’s accelerator, nuclear-physics and detector infrastructures. The project has enabled new experiments with stable, radioactive and neutron beams while advancing accelerator technologies, detectors, beam diagnostics, automation and applications. Its most important achievement is therefore not simply the amount of research access provided, but the way in which science and technology have developed together: new scientific questions have driven technological developments. One of the unique feature of this Transnational Access project was the improvements in  the “services” provided to the users. These included higher intensities and stability of ion beams, improved diagnostics and detectors, to ease of access process. These service improvements in many cases benefited more than one facility.
A broad range of nuclear-physics experiments was enabled through access to complementary European facilities. Using unstable nuclei as beams (that do not exist on earth) at ISOLDE, measurements of single-neutron energies and spectroscopic factors around the doubly magic $^{132}_{50}Sn_{82}$ provided new information on shell structure in neutron-rich nuclei very far from the valley of stability. At GSI/FAIR, R3B experiments investigated nucleon–nucleon correlations in exotic nuclei, while other experiments addressed rare processes such as two-photon nuclear decay. Neutron facilities through measurements of cross sections like neutron capture cross section of radioactive $^{94}$Nb provided data of vital importance to both nuclear astrophysics and energy applications.
EURO-LABS also supported a strong programme using high-intensity stable ion beams, enabling experiments requiring intensities that are not available with radioactive and neutron  beams. At ALTO, stable-beam experiments combined high-resolution $\gamma$-ray detection with fission studies and complementary measurements. At GANIL, inverse-kinematics studies investigated fission mainly through the measurements of isotopic distributions of fission fragments. These provide important inputs to the yet not fully understood dynamics of fission and also for certain applications. At INFN-LNL, AGATA (Advanced Gamma-ray Tracking Array) one of the two arrays in the world, was combined with the PRISMA spectrometer to investigate the collective motion of the nucleons in the nucleus, measure extremely short lifetimes (down to femto secs) and explore the structure in unstable nuclei when both the number of neutrons and protons are identical. At JYFL, stable-beam studies extended precision measurements towards neutron-deficient nuclei and the proton drip line and also for physics beyond the standard model, while GSI/FAIR combined intense heavy-ion beams with storage-ring techniques for precision studies of highly charged ions and rare decay processes. 
CLEAR (Cluster of Low Energy Accelerator Research) is a set of three laboratories  in three different countries which acted as one entity for e.g. have a common Program Advisory Committee, so as to ensure that experiments utilized the most suitable facility.  This is the first time such an endeavour was made possible. A wide number of topic ranging from Accelerator Mass Spectrometry determinations of $^{237}Np$ from seaweed samples characterization of Helium solid targets for experiments, search for dark matter particles in nuclear decays were made at the different CLEAR centres. At IFJ-PAN for e.g., the characterization of the next generation detectors for the R3B (Reactions with Relativistic Radioactive Beams collaboration) for the new $\sim$3 billion euros state-of-the-art FAIR facility in Germany were made. One example of the Transnational Access, highlighting the richness and complementarity of the European teams, was the work done by an  Italian led team  exploiting the complementary expertise and equipment’s in Poland  (IFJ-PAN) and Romania  (IFIN-HH) to address a scientific question. The goal was to measure the temperature and isospin dependence of the oscillations of the neutron skin and also to address the equation of state in neutron star mergers and to understand the isotope abundance in the universe. Various activities illustrate a key EURO-LABS strength: beam development, target technology, advanced detectors and nuclear-physics questions that were developed together, allowing experiments that would otherwise be difficult or impossible and connecting fundamental nuclear research with astrophysics and applications.
The project also demonstrated how the same accelerator infrastructure can address both fundamental science and major technological challenges. At CERN’s HiRadMat (High-Radiation to Materials), intense proton beams were used to investigate radiation damage and thermomechanical effects in materials and accelerator components relevant to future high-power machines. At the same facility, the FIREBALL experiment created laboratory relativistic pair-plasma conditions, bringing accelerator technology into the domain of laboratory astrophysics. The role of nuclei in astrophysical sites was investigated also using Radioactive ion beams. Extreme conditions created at accelerators thus became simultaneously a tool for developing future machines and a laboratory for studying fundamental physics.
One of the strongest examples of the science–technology-health connection is the new FLASH radiotherapy. This is an innovative cancer treatment by the delivery of ultra-high doses. These range from 40 to 100+ Gray/s, of radiation in less than 100 milliseconds. This has to be compared to typical conventional radio therapy of 20 milli gray in a few sec. At GSI/FAIR, EURO-LABS supported work on the technical requirements for ultra-high-dose-rate proton and carbon-ion FLASH, including beam monitoring and dosimetry. At CLEAR (CERN Linear Electron Accelerator for Research), a complementary approach was pursued with Very-High-Energy Electron (VHEE) FLASH, including technologies for producing uniform and conformal irradiation fields and for performing real-time dosimetry at ultra-high dose rates. These activities bring together accelerator physics, beam delivery, detector technology and medical physics, addressing key challenges in translating the FLASH effect towards future clinical applications which improve both the quality and length of the life of cancer patients.
EURO-LABS similarly connected detector development with both fundamental research and applications. European test-beam and irradiation facilities were used to characterize new detector technologies under controlled particle and radiation environments, including Low Gain Avalanche Diodes, ultra-fast silicon detectors and monolithic pixel technologies. The application of fast silicon detectors in proton-therapy beam monitoring, demonstrates how detector technologies developed for particle physics can be adapted to medical applications. Improvements to irradiation systems, beam-profile monitoring, focusing and high-fluence delivery further strengthened Europe’s ability to qualify detectors for demanding future experiments.
A particularly tangible achievement was the development and deployment of an open-source code \textbf{ Geoff } , a Generic Open-Source Framework bringing modern optimization and machine-learning techniques into real accelerator operation. Developed collaboratively between CERN and GSI, Geoff was deployed at both facilities and demonstrated improvement on demanding beam-delivery problems. This code was developed in a versatile manner such that it can be made suitable to any laboratory what ever the instrument/accelerator involved. At CERN, it contributed to increased heavy-ion transmission from the SPS to the LHC. At GSI, for the SIS18 synchrotron it was used to automate optimization reducing injection losses from around 40–45 \% to about 15 \%, reducing tuning time from up to two hours to approximately 15–20 minutes. At the GSI Fragment Separator, Geoff demonstrated automated beam steering and focusing. The work subsequently advanced towards multi-objective and multi-fidelity Bayesian optimization and applications to complex systems such as ion sources and traps. Machine-learning approaches were also investigated for beam-profile classification and the optimization of laser-driven electron sources. This and other related machine learning approaches showed how diagnostics, software, modelling and accelerator control can increasingly form a single optimization loop.
EURO-LABS also contributed to technologies required by the next generation of accelerators. At KARA (Karlsruhe Research Accelerator)), the BESTEX (BEam Screen Test bench Experiment) facility investigated synchrotron-radiation-induced effects on beam-screen materials, including photoelectron production and surface desorption. These measurements provide essential information for controlling electron-cloud effects and vacuum conditions in future high-intensity colliders such as the FCC-ee (Future  Circular electron-electron Collider) at CERN. The latter will be the flagship accelerator for high energy physics  in Europe and in the world. At other centres, accelerator R\&D supported superconducting RF and magnetic technologies, while facilities such as CLEAR also provided platforms for advanced beam diagnostics and novel accelerator applications.
The scientific programme extended beyond nuclear physics and accelerator development into new applications of accelerator technology. Electron-beam irradiation at INCT-RAPID supported research in materials, chemistry, biology and environmental applications, while CLEAR combined electron-beam technology with FLASH, advanced diagnostics and detector studies. These activities demonstrate how specialised accelerator infrastructure can generate scientific results well beyond the traditional boundaries of accelerator and nuclear physics.
Another important achievement was Theo4Exp, which brings theoretical calculations closer to the experiments. By providing common access to nuclear-theory tools, EURO-LABS enabled calculations to support the planning and interpretation of experiments both for nuclear structure and nuclear reactions. The users of this virtual infrastructure are supported by experts some of them who have themselves written these state-of-the-art codes.  EURO-LABS has also established a common Open Science ecosystem across its participating facilities, significantly improving the FAIRness, accessibility and reuse of nuclear-physics data. A key achievement is OpenNP, a FAIR-compliant service (Findability, Accessibility, Interoperability, and Reuse of digital assets) providing a catalogue of nuclear-physics datasets and linking them with associated software and publications. In parallel, EURO-LABS has developed a federated authentication and authorisation infrastructure (AAI), enabling secure and streamlined access to shared services across participating laboratories. This infrastructure is already supporting services such as OpenNP and Theo4EXP and provides a common basis for future cross-facility developments. A common Data Management Plan provides the framework for managing project and Transnational Access data through standardized metadata, persistent identifiers, appropriate repositories and licensing, while addressing security, privacy and intellectual property requirements. Together with solutions for remote access, data sharing and monitoring, and dedicated training activities, these developments have created sustainable infrastructure and practices that will extend beyond EURO-LABS. They strengthen reproducibility, cross-facility collaboration and the integration of nuclear-physics research into the wider European Open Science ecosystem, including EOSC (European Open Science Cloud).
One of the most successful and distinctive aspects of EURO-LABS was its hands-on training programme, training by doing: building the next generation of experts. This was designed to give young scientists and engineers direct experience with the research infrastructures, technologies and experimental methods used across the research infrastructures. Rather than relying on lectures, the schools mainly put trainees in front of real accelerators, detectors, control systems and experimental setups, allowing them to operate equipment, take data, solve technical problems and work alongside experienced researchers. The programme included basic and advanced Training Schools (four and five respectively), as well as specialised InTraNS -Instrumentation and Training for Nuclear Spectroscopy and Reaction Dynamics -  (6 schools and workshops). The first basic School at IFIN-HH in was built around training in the laboratory and experiments on the tandem accelerators for a 12-day programme..  At CERN, trainees worked directly with CLEAR, ISOLDE, AD and the PS Booster, gaining experience in accelerator control, beam characterisation, steering and superconducting-cavity operation. In parallel, InTraNS provided specialised practical training in various areas ranging from HPGe detector operation, testing and repair, analysis to data analysis using AGATA. In 2025, the CERN accelerator school brought together 19 early-career scientists and engineers from 10 countries for an intensive programme combining lectures, simulations and practical beam operations. Subsequent schools extended the experience to accelerator applications, accelerator operation, detector technology, nuclear spectroscopy and open science. The exceptionally positive feedback from participants highlighted the value of direct contact with equipment, the diversity of facilities and technologies, and interaction with both peers and experienced specialists. EURO-LABS therefore did more than providing access to infrastructures: it transferred practical know-how across laboratories and generations, helping to build the scientific and technical community needed to exploit Europe’s accelerator facilities in the future. To optimize the use of funds, a limited number of students were supported to already existing schools (5 partially supported and 6 co-sponsored). Figure 2 shows the participants in a few of hands-on schools organized by EURO-LABS. In addition to this ECT*, the European Centre for Theoretical Studies in Nuclear Physics and related areas conducted a large number and a variety of  5-day workshops with world leading lecturers. This included a dedicated hands-on workshop for exploiting  Theo4EXP facilities.
Taken together, the part of the achievements discussed here and the rest which can be seen in the various reports and presentations at the various annual meetings, international conferences and publications, show the breadth of EURO-LABS’ impact. These reports  can be found at PUT EURO-LABS WEBSITE
%PUT EURO_LABS WEBSITE%.
The project has enabled new measurements of exotic nuclei and nuclear reactions, created opportunities for extreme-beam and laboratory-astrophysics experiments, and supported technologies for future accelerators, detectors and medical applications. It has introduced machine learning and automated optimisation into accelerator operation, strengthened the connection between theory and experiment, improved access to advanced infrastructures and created new approaches to training and scientific data.
The distinctive legacy of EURO-LABS is the building of a common team of the very large and diverse “super” sub-atomic community of Europe. It provided effective access to a network of 44 Research Infrastructures from 15 countries (including 3 RIs with Virtual Access) to conduct curiosity-based research, addressing fundamental questions and technological challenges and advancing projects with broad societal impact. It fostered knowledge sharing between scientific fields and enhancing Europe’s potential for successfully facing future challenges. Through this large network, EURO-LABS ensured the diversity and active support of researchers from different nationalities, gender, age, grade, and with a variety of professional expertise. EURO-LABS has not only increased access to accelerator-based science—it has expanded what these facilities can do and created new routes from fundamental research to advanced technology and built a foundation for a new access policy to Research Infrastructures of Europe. The community looks forward to working together in even more diverse projects of the European Union and thus contribute to the excellence of European and world Science. 


\begin{figure}[htb]
   \centering
   \includegraphics[width=\columnwidth]{MOP7145_f2.pdf}
    \caption{Age-gender distribution for the EURO-LABS supported users in TNA projects.}
   \label{fig:age-gender}
   % \vspace{-2mm}
\end{figure}
Particular attention was devoted to diversity and inclusiveness in the allocation of support, considering aspects such as gender balance, nationality, institutional origin, and career stage of the supported users, as illustrated in Figs.~\ref{fig:age-gender}, \ref{fig:home-inst}, and \ref{fig:nationality}.
\begin{figure}[htb]
   \centering
   \includegraphics[width=\columnwidth]{MOP7145_f3.pdf}
    \caption{Home Institute origin for the EURO-LABS supported users in TNA projects.}
   \label{fig:home-inst}
\end{figure}

\begin{figure}[htb]
   \centering
   \includegraphics[width=\columnwidth]{MOP7145_f4.pdf}
    \caption{Researcher Nationality for the EURO-LABS supported users in TNA projects.}
   \label{fig:nationality}
%   \vspace{-5mm}
\end{figure}

The EURO-LABS Newsletter~\cite{newsletter} published periodically collects scientific and community highlights of the project. Key scientific highlights include a report on the FLASH effect radiotherapy R\&D aimed at destroying or controlling tumor with Ultra-high dose rate (UHDR; >\SI{40}{Gy/s} irradiation with minimal normal tissue toxicity, realized in GSI-FAIR and CERN-CLEAR facilities of EURO-LABS. Or, the series of experiments at CERN-HiRadMat aimed at producing a lab-analogue of the jets of matter and antimatter that stream out of some black holes and neutron stars. In addition, the release of the first version of the \textbf{openNP catalogue}~\cite{opennp}, one of the main objectives of the project. OpenNP is a web-based service offering a centralized and FAIR-compliant metadata catalog for nuclear physics datasets, as those produced through EURO-LABS TNA projects.

The training activities within EURO-LABS constitute an important pillar of the project. This objective was successfully pursued through the organisation of nine training schools across Europe, covering both introductory and advanced-level topics. A particular highlight are the three Advanced Hands-On Training Schools on Operation of Accelerators (\textbf{ATSOA}) held at CERN in 2024-2026, which provided selected students with a unique practical training experience in accelerator operation in a wide spectrum of accelerators~\cite{atsoa-paper}.

\subsection{Summary}
The most durable legacy of EURO-LABS is therefore distributed. It is present in scientific results, but also in improved infrastructure services, software, detector technologies, experimental methods, training networks, open-data practices and—perhaps most importantly—in relationships between people and facilities.
Several developments are readily identifiable as reusable assets: the Geoff optimization framework; Theo4Exp and its interoperable theory nodes; OpenNP and associated FAIR-oriented services; new beam diagnostics and monitoring systems; detector test and irradiation capabilities; improved materials-testing services; and practical training formats that can be repeated by the participating communities.
The human legacy is harder to measure but is supported by the D5.2 statistics and the training programme. Thousands of unique users were brought into a shared European access environment, with strong representation of early-career scientists and participation extending well beyond the countries hosting the facilities. The project also created repeated opportunities for engineers, technicians and researchers to learn from one another.
The four-year trajectory from KOM to FINE is therefore best understood as a transition from coordination to community. The project did not remove the differences between facilities or scientific disciplines; instead, it made those differences productive by creating mechanisms through which they could interact.
EURO-LABS demonstrated a model in which Transnational Access, scientific collaboration, technical development, training, open science and diversity reinforce one another. The project provided substantial access to European Research Infrastructures, supported a broad range of experiments and technology programmes, and used the requirements of users to drive concrete improvements in services.
Its scientific breadth is visible in studies of nuclei far from stability, neutron-induced reactions, fission, detector technologies, accelerator materials, plasma and electron beams, FLASH-related developments and future accelerator technologies. Its technical legacy is visible in improved diagnostics, controls, irradiation and testing capabilities, open-source optimization, theory services and data infrastructure. Its community legacy is visible in the geographical, gender and age diversity of the users, and in the repeated hands-on opportunities offered to the next generation.
The lasting lesson is that research infrastructures become more powerful when access is accompanied by mechanisms for people, ideas and technologies to move between them. EURO-LABS created such mechanisms across nuclear physics, accelerator science and detector development. The four years therefore leave not a single flagship result, but a connected set of scientific, technological and human capabilities that can continue to be used and developed.

%EURO-LABS has successfully established a unique collaborative framework that, for the first time within a single European %initiative, brings together the Nuclear Physics, High-Energy Physics, and Accelerator R\&D communities around a common %ecosystem of research infrastructures. The project enabled efficient transnational access to leading facilities across %Europe and supported a broad range of scientific activities through its TNA programme. EURO-LABS provided scientists, in %particular researchers at early stages of their careers, with the opportunity to realise high-quality projects fostering %novel research and technological developments. In addition to its scientific output, the project contributed significantly %to training, community integration, and the development of advanced services and methodologies, demonstrating the value of a %coordinated cross-disciplinary RI approach at the European level.

%
% only for "biblatex"
%
% JACoW-paper template trailer (bibliography, \end{document}) commented out: it broke the report build
%\ifboolexpr{bool{jacowbiblatex}}%
%	{\printbibliography}%
%	{%
%	% "biblatex" is not used, go the "manual" way
%	
%	%\begin{thebibliography}{99}   % Use for  10-99  references
%	\begin{thebibliography}{9} % Use for 1-9 references
%	
%%	\bibitem{euro-labs}
%%		European Laboaratories for Accelerator Based Science - HORIZON-INFRA-2021 — SERV-01-07,
%%		\url{https://web.infn.it/EURO-LABS/}
%
%	%\bibitem{euro-labs}
%%European Laboratories for Accelerator Based Science,
%%\url{https://web.infn.it/EURO-LABS/}
%	
%%	\bibitem{horizon-europe}
%%		\textit{Horizon Europe - The EU Research \& Innovation Programme 2021-2027.},
%%        \url{https://research-and-innovation.ec.europa.eu/funding/funding-opportunities/funding-programmes-and-open-calls/horizon-europe_en}
%
%%	\bibitem{horizon-europe}
%%	Horizon Europe - The EU Research \& Innovation Programme 2021-2027, \\
%%	\url{https://research-and-innovation.ec.europa.eu/funding/funding-opportunities/funding-programmes-and-open-calls/%horizon-europe_en}
%
%%    \bibitem{geoff}
%%        P.~Madysa, S.~Appel, V.~Kain, and M.~Schenk,
%%        "Geoff: The generic optimization framework \& frontend for particle accelerator controls",
%%        \textit{SoftwareX}, vol. 32, p. 102335, Dec. 2025.
%%        \doi{10.1016/j.softx.2025.102335}
%
%%    %\bibitem{eurolabs-web}
%%    %EURO-LABS, 
%%    %\url{https://web.infn.it/EURO-LABS/}
%    
%%    \bibitem{newsletter}
%%        M.~Colonna,
%%        \textit{EURO-LABS newsletter}, \\
%%        \url{https://web.infn.it/EURO-LABS/category/media/newsletter/}
%
%%    \bibitem{newsletter}
%%	EURO-LABS newsletter, \\
%%	\url{https://web.infn.it/EURO-LABS/category/media/newsletter/}
%
%%      \bibitem{opennp}
%%        A.~Matta, A.~Lemasson, and M.~J.~G.~Borge,
%%        "D5.3 Release of the first functional version of the openNP catalog", 
%%        EURO-LABS - EUROpean Laboratories for Accelerator Based Science, Rep. EURO-LABS-Del-5-3, 2026. 
%%        \doi{10.5281/zenodo.17020290}
%        
%%    \bibitem{atsoa-paper}
%%        F.~Asvesta \emph{et al.},
%%        "ATSOA at CERN: A Hands On Accelerator Course in the EURO-LABS Framework",
%%        presented at IPAC'26,
%%        Deauville, France, May 2026, paper TUP8014, this conference.
%
%%	\end{thebibliography}
%
%  
%%} % end \ifboolexpr
%%
%% for use as JACoW template the inclusion of the ANNEX parts have been commented out
%% to generate the complete documentation please remove the "%" of the next two commands
%% 
%%%%\newpage
%
%%%%\include{annexes-A4}
%
%\end{document}

